\documentclass[11pt,a4paper]{article}

% --- Encodage / langue ---
\usepackage[utf8]{inputenc}
\usepackage[T1]{fontenc}
\usepackage[french]{babel}

% --- Mise en page ---
\usepackage[margin=2.2cm]{geometry}
\usepackage{booktabs}
\usepackage{array}
\usepackage{longtable}
\usepackage{graphicx}
\usepackage{amsmath,amssymb}
\usepackage{xcolor}
\usepackage{enumitem}
\usepackage{fancyhdr}
\usepackage[colorlinks=true,linkcolor=black,urlcolor=blue,citecolor=black]{hyperref}

\setlist{nosep,leftmargin=1.4em}
\renewcommand{\arraystretch}{1.15}

% En-tete / pied de page
\pagestyle{fancy}
\fancyhf{}
\lhead{\small Recommender Gaming}
\rhead{\small M.\,Y.\,Cheikh Med Vall — 25\,239}
\cfoot{\small\thepage}
\renewcommand{\headrulewidth}{0.3pt}

% Boite « capture a inserer »
\newcommand{\captureplaceholder}[1]{%
  \begin{center}
  \fbox{\parbox[c][3.2cm][c]{0.82\linewidth}{\centering\itshape\color{gray} Capture à insérer — #1}}
  \end{center}}

\begin{document}

% ======================= TITRE =======================
\begin{center}
  {\LARGE\bfseries Système de recommandation de jeux vidéo}\\[2pt]
  {\large Hybride pondéré contenu + collaboratif, avec front-end LLM et voix locale}\\[10pt]
  {\normalsize Projet — Systèmes de recommandation \quad\textbullet\quad SupNum}\\[2pt]
  {\normalsize \textbf{Mouhamedou Yahya Cheikh Med Vall} — Matricule 25\,239}\\[2pt]
  {\normalsize Sujet imposé : \emph{Recommender gaming} \quad\textbullet\quad Juin 2026}
\end{center}

\vspace{0.4em}
\hrule
\vspace{1em}

% ======================= 1. PROBLEME =======================
\section{Problème}

L'objectif est de recommander des jeux vidéo à partir d'une requête en langage
naturel, tapée ou dictée, et de restituer les résultats à l'écran et à voix
haute. Le c\oe ur du système est un \textbf{hybride pondéré} au sens de Burke :
il combine deux façons indépendantes de mesurer la similarité entre deux jeux —
par leurs \emph{métadonnées} (contenu) et par \emph{qui les joue} (collaboratif)
— puis recommande les plus proches voisins d'un jeu d'amorçage (\emph{seed}).

Le projet est volontairement \emph{académique et pédagogique} : la clarté et la
correction priment sur l'échelle ou le vernis industriel. Chaque composant doit
rester assez simple pour être expliqué en soutenance.

Conformément à l'énoncé, le sujet « Recommender gaming » s'appuie sur un jeu de
données de recommandation validé par l'enseignante. L'application livrée est
développée en \textbf{React} (et non Streamlit), modification validée au
préalable avec l'enseignante ; le reste des livrables (notebook de
prétraitement, README, évaluation) est respecté.

% ======================= 2. DATASET =======================
\section{Jeu de données}

Le jeu de données est « Game Recommendations on Steam » (Kaggle). Il est retenu
parce qu'il porte \textbf{les deux signaux} dont l'hybride a besoin : des
métadonnées de jeux pour le bras contenu, et des interactions utilisateur–jeu
pour le bras collaboratif. Quatre fichiers bruts :

\begin{itemize}
  \item \texttt{games.csv} — 50\,872 jeux $\times$ 13 colonnes (titre, prix,
        ratio d'avis positifs, date, plateformes).
  \item \texttt{games\_metadata.json} — 50\,872 lignes (\texttt{app\_id},
        \texttt{description}, \texttt{tags}) ; 49\,628 ont $\geq 1$ tag,
        40\,499 ont une description non vide. \emph{Carburant du bras contenu.}
  \item \texttt{recommendations.csv} — $\approx$ 41 millions de lignes d'avis
        (présence d'une review = interaction). \emph{Carburant du bras
        collaboratif.}
  \item \texttt{users.csv} — table des utilisateurs (optionnelle dans le design).
\end{itemize}

\paragraph{Échantillonnage (choix justifié).} L'ensemble complet des
interactions est trop lourd pour un calcul item--item sur portable. On retient
donc le \textbf{catalogue des 6\,000 jeux les plus interagis} (et non un seuil
arbitraire de tags : ce seuil avait initialement évincé des titres phares comme
Witcher~3, GTA~V ou CS:GO, ce qui a été corrigé). On échantillonne ensuite
\textbf{150\,000 utilisateurs} (graine \texttt{42}) ayant au moins 5
interactions dans le catalogue. La matrice d'interactions finale est de taille
$150\,000 \times 6\,000$, avec une \textbf{sparsité $\approx 99{,}81\,\%$}
(1\,680\,531 valeurs non nulles).

% ======================= 3. FEATURES =======================
\section{Caractéristiques (features)}

\paragraph{Bras contenu.} Chaque jeu est représenté par une « soupe » de texte :
ses tags (espaces internes remplacés par des \texttt{\_}, bloc répété 3 fois pour
dominer la prose) concaténés à sa description nettoyée. Une liste d'arrêt
(\emph{stoplist}) retire $\approx$ 35 tokens de « plomberie » plateforme
(\texttt{single\_player}, \texttt{steam\_achievements}, \dots) qui décrivent
l'infrastructure plutôt que le contenu. Cette soupe est vectorisée en
\textbf{TF-IDF} : matrice $6\,000 \times 8\,681$ (8\,681 termes de vocabulaire),
lignes normalisées L2 — le produit scalaire donne donc directement un cosinus
dans $[0,1]$.

\paragraph{Bras collaboratif.} Chaque jeu est un vecteur sur l'espace des
utilisateurs (signal \emph{binaire} : présence d'une review). La similarité est
un \textbf{cosinus item--item pondéré par IDF}, la pondération neutralisant
l'effondrement vers les jeux les plus populaires.

\paragraph{Filtres.} Les contraintes extraites de la requête (\texttt{max\_price},
\texttt{tags}, \texttt{genres}) ne servent pas à la similarité : elles
\emph{contraignent} la liste de résultats après le score, jamais avant.

% ======================= 4. METHODE =======================
\section{Méthode}

\subsection{Hybride pondéré}

Le score de similarité entre deux jeux $a$ et $b$ est une combinaison convexe
des deux bras, réglée par un unique poids $\alpha$ :
\[
  \mathrm{score}(a,b) \;=\; \alpha \cdot \mathrm{sim}_{\text{contenu}}(a,b)
  \;+\; (1-\alpha)\cdot \mathrm{sim}_{\text{collab}}(a,b).
\]
Les deux bras étant sur des échelles de magnitude différentes, chaque vecteur de
scores est \textbf{re-normalisé par son maximum} avant le mélange — sans cette
étape, le mélange « bascule » d'un bras à l'autre au lieu d'interpoler. Le poids
par défaut retenu est $\alpha = \mathbf{0{,}25}$ (provisoire ; justification en
§\,\ref{sec:eval}).

\subsection{Routage et démarrage à froid (cold-start)}

Tout jeu n'a pas les deux signaux. Le système \emph{route} donc chaque amorce :
\begin{itemize}
  \item vecteur contenu nul $\rightarrow$ \texttt{cf\_only} (5 jeux : des MMO F2P
        sans métadonnées exploitables) ;
  \item signal collaboratif faible, $< 50$ interactions $\rightarrow$
        \texttt{content\_only} (2\,068 jeux, $\approx 34\,\%$ du catalogue) ;
  \item les deux présents $\rightarrow$ \texttt{blend} (mélange pondéré) ;
  \item aucun $\rightarrow$ repli \texttt{popularity}.
\end{itemize}

\subsection{Compréhension de la requête (LLM)}

Un modèle local \texttt{llama3.1:8b} (via Ollama) traduit la requête libre en une
structure \texttt{\{ seed, filters \}}. Il est utilisé \textbf{uniquement pour
analyser} la requête (et, le cas échéant, proposer un jeu d'amorce
représentatif) ; il \emph{n'écrit jamais} les recommandations — celles-ci
viennent toujours de l'hybride. Le titre proposé est résolu vers le catalogue
par un appariement \texttt{max(ratio\_caractères, recouvrement\_tokens)} avec un
seuil d'admission de \textbf{0{,}80} : en dessous, le système répond honnêtement
« je n'ai pas reconnu de jeu » plutôt que de s'accrocher à un mauvais voisin.

\subsection{Amorçage par ambiance (vibe-seeding)}

Quand la requête décrit une \emph{ambiance} sans nommer de jeu (« quelque chose
de relaxant et atmosphérique »), le système dérive néanmoins une amorce : le LLM
propose le jeu le plus représentatif de l'ambiance, ce titre est \emph{ancré}
par le résolveur (anti-hallucination), et à défaut un repli TF-IDF transforme le
texte dans l'espace contenu pour en chercher le plus proche voisin. Un seuil de
confiance garantit le même principe : amorcer avec assurance ou décliner avec
assurance.

\subsection{Interface et voix}

Le front-end React capture la requête (texte ou micro) et affiche les résultats ;
il \textbf{ne calcule rien} (aucun tri, score ou filtre côté client). La voix est
servie \emph{localement} par le back-end : transcription par \texttt{faster-whisper}
(\texttt{POST /stt}) et synthèse par \texttt{Piper} (\texttt{POST /tts}). Ce choix
local supprime la dépendance au cloud, corrige la mauvaise qualité de la synthèse
navigateur et fonctionne au-delà de Chrome.

\subsection{Flux d'une requête}

\begin{enumerate}
  \item \textbf{Requête} — tapée, ou dictée puis transcrite par \texttt{/stt}.
  \item \textbf{Analyse LLM} — la requête devient \texttt{\{ seed, filters \}} ;
        si \texttt{seed} est nul, l'amorçage par ambiance prend le relais.
  \item \textbf{Score hybride} — cosinus contenu + cosinus CF, mélangés par $\alpha$.
  \item \textbf{Filtrage + tri} — application des filtres, Top-$N$ voisins de l'amorce.
  \item \textbf{Sortie} — affichage React + lecture vocale via \texttt{/tts}.
\end{enumerate}

% ======================= 5. CAPTURES =======================
\section{Captures}

\captureplaceholder{page d'accueil : barre de requête (texte + micro) et résultats}
\captureplaceholder{requête vocale « a hard roguelike like Hades » et voisins recommandés}
\captureplaceholder{réponse \emph{no\_seed} guidée + tier de confirmation « did you mean\dots ? »}

% Pour insérer une vraie capture, remplacer la commande \captureplaceholder par :
% \begin{figure}[h]\centering\includegraphics[width=0.82\linewidth]{figs/accueil.png}
% \caption{...}\end{figure}

% ======================= 6. EVALUATION =======================
\section{Évaluation}\label{sec:eval}

\paragraph{Protocole.} \emph{Leave-one-out} sur les interactions : pour 1\,000
utilisateurs échantillonnés (graine \texttt{42}), on retire une interaction et on
mesure si elle est retrouvée dans le Top-$k$ des voisins. Le score d'un
utilisateur est calculé sur \textbf{l'ensemble de son profil} (et non sur
quelques amorces) : on a vérifié que ce protocole « profil complet » relève le
taux de réussite d'$\approx 35\,\%$ en relatif par rapport à un raccourci à 3
amorces (0{,}119 $\rightarrow$ 0{,}161 à $\alpha=0$), donc on l'adopte comme
protocole honnête. On balaie $\alpha$ de 0 (collaboratif pur) à 1 (contenu pur).

\paragraph{Métrique.} Avec une seule cible retirée par essai, la précision@$k$
est mécaniquement plafonnée à $1/k$ ; on met donc en avant le \textbf{taux de
réussite (hit-rate@$k$)}, égal au rappel@$k$ ici. La précision@10 vaut
exactement hit-rate@10$/10$.

\paragraph{Ablation — balayage de $\alpha$} (profil complet, $n=1\,000$,
graine \texttt{42}) :

\begin{center}
\begin{tabular}{lccccc}
\toprule
$\alpha$ & 0{,}00 & \textbf{0{,}25} & 0{,}50 & 0{,}75 & 1{,}00 \\
\midrule
Hit-rate@10   & 0{,}161 & \textbf{0{,}168} & 0{,}148 & 0{,}071 & 0{,}029 \\
Hit-rate@20   & 0{,}247 & \textbf{0{,}250} & 0{,}205 & 0{,}103 & 0{,}046 \\
Couverture@10 & 0{,}201 & 0{,}247 & 0{,}305 & 0{,}368 & \textbf{0{,}387} \\
Diversité@10  & 0{,}933 & 0{,}908 & 0{,}865 & 0{,}814 & 0{,}793 \\
\bottomrule
\end{tabular}
\end{center}

\paragraph{Lecture honnête — un gain sans coût.} À $\alpha=0{,}25$, le taux de
réussite (\textbf{0{,}168}) est \emph{statistiquement indiscernable} du
collaboratif pur ($\alpha=0$ : 0{,}161 ; écart $+0{,}007$, erreur-type
$\approx 0{,}012$) : l'exactitude est \textbf{maintenue}. Dans le même temps la
\textbf{couverture catalogue augmente} de $+4{,}6$ points (0{,}201 $\rightarrow$
0{,}247). Le mélange n'est donc pas un compromis perte-de-précision-contre-
couverture : c'est une amélioration de \textbf{Pareto} — plus de couverture pour
aucun coût d'exactitude mesurable. C'est ce qui justifie $\alpha=0{,}25$. (Au-delà
de $0{,}5$, l'exactitude chute nettement : le mélange équilibré est le bon réglage.)

\paragraph{Sur la diversité (mise en garde).} La diversité décroît avec $\alpha$,
mais cette mesure est \textbf{circulaire} : elle évalue une liste avec le cosinus
contenu, alors que le bras contenu a contribué à produire cette liste. On
s'appuie donc sur la \textbf{couverture} comme signal non circulaire, et la
diversité n'est citée qu'avec cette réserve.

\paragraph{Pourquoi la couverture est structurellement réelle.} Les angles morts
des deux bras sont \emph{disjoints} : le bras contenu est aveugle sur 5 jeux
(vecteur nul), que le collaboratif couvre ; le bras collaboratif est à faible
signal sur 2\,068 jeux, que le contenu couvre. Chaque bras comble l'angle mort de
l'autre — la couverture mesurée ci-dessus n'est pas un artefact, c'est la raison
d'être de l'hybride. Tous les chiffres proviennent du notebook
\texttt{notebooks/eval.ipynb} et de \texttt{artifacts/eval\_metrics.json}.

% ======================= 7. LIMITES =======================
\section{Limites}

\begin{itemize}
  \item \textbf{Le taux de réussite brut reste modeste} (hit-rate@10 $\approx$
        17\,\%) : c'est attendu pour un leave-one-out à une seule cible positive
        sur un catalogue de 6\,000 items ; la métrique mesure le rappel d'un
        co-jeu, pas la pertinence subjective.
  \item \textbf{L'hybride se défend sur la couverture} : à $\alpha=0{,}25$
        l'exactitude est maintenue (indiscernable du collaboratif pur) et la
        couverture gagne $+4{,}6$ points — honnêteté assumée plutôt qu'un
        graphique flatteur.
  \item \textbf{La diversité mesurée est circulaire} (cosinus contenu jugeant une
        liste façonnée par le bras contenu) : on s'appuie sur la couverture.
  \item \textbf{5 jeux sans contenu} (MMO F2P dont l'API Steam ne renvoie rien) :
        résidu accepté, traités par le repli collaboratif.
  \item \textbf{TF-IDF est lexical} : l'amorçage par ambiance par TF-IDF seul
        confond « relaxant » avec la présence littérale du mot ; d'où le recours
        au LLM (ancré par le résolveur) pour la sémantique des ambiances.
  \item \textbf{Échelle réduite} : 6\,000 jeux et 150\,000 utilisateurs
        échantillonnés pour rester traitable sur portable ; le système n'est pas
        dimensionné pour le catalogue Steam complet.
  \item \textbf{Latence vocale} : l'analyse LLM domine le temps de réponse
        ($\approx 4$\,s à chaud) ; un état de chargement masque l'attente, mais
        la chaîne STT $\rightarrow$ LLM $\rightarrow$ reco $\rightarrow$ TTS reste
        séquentielle.
\end{itemize}

% ======================= 8. EXECUTION =======================
\section{Installation et exécution}

\begin{verbatim}
# Back-end (FastAPI + artefacts pré-calculés)
conda activate ds
pip install -r requirements.txt
uvicorn src.api:app --port 8765

# LLM local (dans un terminal séparé)
ollama serve            # modèle requis : llama3.1:8b

# Front-end (Vite + React)
cd web
npm install
npm run dev             # http://localhost:5173
\end{verbatim}

\noindent Déterminisme : toutes les étapes d'échantillonnage utilisent la graine
\texttt{42} ; les artefacts (\texttt{games.json}, matrices de similarité,
vectoriseur TF-IDF) sont régénérables depuis \texttt{src/} et ne sont jamais
édités à la main.

% ======================= CONCLUSION =======================
\section*{Conclusion}
\addcontentsline{toc}{section}{Conclusion}

Le livrable est un recommandeur de jeux \emph{hybride pondéré} fonctionnel et
explicable : deux bras cosinus mélangés par un seul $\alpha$, un front-end LLM
local qui transforme une requête libre (ou une simple ambiance) en amorce sans
jamais inventer les recommandations, et une voix entièrement locale. L'évaluation
assume un résultat nuancé et le mesure : à $\alpha=0{,}25$ l'exactitude est
maintenue (indiscernable du collaboratif pur) tandis que la couverture du
catalogue augmente — une amélioration de Pareto, et non un score unique gonflé.
C'est cette honnêteté, plutôt qu'un chiffre flatteur, qui correspond le mieux à
l'objet réel d'un système de recommandation.

\end{document}
